% !TeX program = pdflatex
% Gallery starter; see TEMPLATE-README.txt for attribution and licenses.
\documentclass[11pt,a4paper]{article}
\usepackage[T1]{fontenc}
\usepackage{lmodern,microtype}
\usepackage[margin=24mm]{geometry}
\usepackage{graphicx,booktabs,tabularx,enumitem,amsmath,amssymb,xcolor,titlesec,fancyhdr}
\usepackage[hidelinks]{hyperref}
\definecolor{accent}{HTML}{B45C25}\definecolor{quiet}{HTML}{526171}
\setlength{\parindent}{0pt}\setlength{\parskip}{5pt}
\setlist[itemize]{leftmargin=*,itemsep=3pt}
\titleformat{\section}{\Large\sffamily\bfseries\color{accent}}{\thesection}{0.6em}{}
\titleformat{\subsection}{\sffamily\bfseries}{\thesubsection}{0.6em}{}
\titlespacing*{\section}{0pt}{12pt}{5pt}
\titlespacing*{\subsection}{0pt}{10pt}{4pt}
\pagestyle{fancy}\fancyhf{}\fancyhead[L]{\small\sffamily Laboratory Report}\fancyhead[R]{\small\sffamily Alex Morgan}\fancyfoot[R]{\small\thepage}
\renewcommand{\headrulewidth}{0pt}
\newcommand{\eyebrow}[1]{{\small\sffamily\bfseries\color{accent}#1}\par\vspace{5mm}}
\newcommand{\hero}[2]{{\fontsize{32}{36}\selectfont\sffamily\bfseries #1\par}\vspace{4mm}{\large\color{quiet}#2\par}\vspace{6mm}{\color{accent}\rule{\linewidth}{1.5pt}}\vspace{5mm}}

\begin{document}
\eyebrow{ACADEMIC / 2026}
\hero{Measuring with confidence}{Experiment 04 / Physics Laboratory}
\begin{tabularx}{\linewidth}{XX}\textbf{Researcher} Alex Morgan & \textbf{Date} 1 November 2026\\\textbf{Partner} Jordan Lee & \textbf{Supervisor} Dr Taylor\\\end{tabularx}
\section{Objective}Estimate the spring constant from repeated measurements of displacement under a known load.
\section{Equipment and procedure}
\begin{enumerate}\item Mount the spring and record its unloaded length.\item Add a known mass and wait for the system to settle.\item Measure displacement three times for each load.\item Plot force against displacement and estimate the gradient.\end{enumerate}
\section{Observations}
\begin{center}\begin{tabular}{rrrr}\toprule Mass (kg) & Force (N) & Displacement (m) & Uncertainty (m)\\\midrule 0.10 & 0.98 & 0.020 & 0.001\\0.20 & 1.96 & 0.039 & 0.001\\0.30 & 2.94 & 0.060 & 0.001\\\bottomrule\end{tabular}\end{center}
\section{Analysis}Using Hooke's law, $F=kx$, the approximate spring constant is $k=49\,\mathrm{N\,m^{-1}}$. These values are illustrative.
\section{Discussion}Discuss systematic error, measurement precision, and whether the model is appropriate over the full range of observations.
\end{document}
